\documentclass[10pt]{article}
\usepackage{tmlr}
\usepackage{amsmath,amssymb}
\usepackage{booktabs,longtable,array,calc}
\usepackage{graphicx}
\usepackage{xcolor}
\usepackage{microtype}
\usepackage{hyperref}
\usepackage{url}
\usepackage{iftex}
\providecommand{\tightlist}{%
  \setlength{\itemsep}{0pt}\setlength{\parskip}{0pt}}
\title{Typed and Scoped Partial-Knowledge State for Research Decisions}
% Keep the source itself anonymous; tmlr.sty also hides authors under review.
\author{Anonymous Authors}
\date{}

\begin{document}
\maketitle

\begin{abstract}
Research agents act not only on facts about the world but on partial
knowledge about those facts: which unknowns belong to which interface
class, which failure records remain applicable, which uncertainties can
change a decision, and which certificates survive representation
changes. We test a bounded mechanism question: \textbf{when competing
methods receive the same underlying serialized information, can explicit
type and scope change the downstream decision?} Across six prospectively
frozen exact-synthetic studies, we isolate unknown-feasibility probing,
failure-receipt reopening, interval verification, evidence transport,
next-experiment selection and representation reminting. Typed-prior
value-of-information probing reaches mean utility 3.291 versus 2.180 for
the identical planner with a flattened prior; scope-bound failure
reopening reaches 3.199 pooled utility while unscoped change reopening
falls to -7.813; decision-boundary-targeted verification reduces
scalarized regret to 0.1096 versus 0.2518 for random verification at the
same budget; full-chain transport checking detects all 200 registered
laundering chains at zero false positives in the frozen synthetic model;
decision-coupled experiment selection reaches utility 9.266 versus 7.121
for pure information gain while avoiding registered entropy decoys; and
typed remint/transport reaches 9.421 in the mixed-transport regime
versus 7.157 for matched re-derivation while tying exactly when
reminting is unnecessary. Two first-right-of-refusal negatives bound the
claim: an ideal VOI donor closes the allocation-policy residual in a
separate typed-failure-state study, and a model-selection donor absorbs
the original crossover-prediction residual. We explicitly subtract
current prior art: STALE \citep{chao2026stale} studies agent-memory
invalidation, ContextNest \citep{sulpovar2026contextnest} studies
governed/versioned/provenanced context, and
value-of-information/provenance methods are established donors. Our
narrower result is exact-synthetic \textbf{mechanism isolation at
matched information}: in registered worlds where a downstream
responsibility depends on a distinction carried by type/scope, erasing
that distinction changes the decision. No result establishes real-agent
safety, scientific productivity, cryptographic security or deployment
performance.
\end{abstract}

\hypertarget{introduction}{%
\subsection{Introduction}\label{introduction}}

A long-running research process accumulates more than observations. It
accumulates state about its own knowledge: an interface is unknown but
belongs to a particular class; a prior attempt failed under a named
representation and access contract; a cost is an interval rather than a
point; a certificate was transported through multiple edits; or a
representation changed after a failure/certificate was recorded.

Many underlying ingredients are mature. Value-of-information methods
\citep{russell1991rightthing} ask which measurement is worth purchasing.
Database and workflow provenance
\citep{buneman2001provenance, cheney2009provenance} track how results
were derived. Agent memory systems increasingly detect and revise stale
state. Context-governance systems track versions, provenance, selectors
and artifact identity. ORION-04 does not claim any of these primitives
as new.

Two current donors are especially important.

\textbf{STALE} asks whether an agent can recognize that remembered state
is no longer valid after later evidence creates an implicit conflict,
and whether that revision propagates into downstream behavior. That is a
genuine \emph{memory invalidation} problem.

\textbf{ContextNest} instead treats provenance, version identity and
deterministic context selection as governance infrastructure beneath
retrieval. A relevant artifact can be stale, unapproved or impossible to
reconstruct even when retrieval itself works.

ORION-04 studies a narrower object than either. We assume the compared
non-oracle policies receive the same serialized facts and ask whether
explicit \textbf{type/scope fields change a registered downstream
responsibility}. In N4-B, for example, the old failure need not become
false because later evidence contradicts it; it becomes potentially
inapplicable because one of the exact context coordinates on which that
receipt depended changed. We call this \textbf{scope invalidation}, not
stale-memory detection.

The same distinction recurs across the other worlds. A high-entropy
unknown can be decision-irrelevant. A certificate can be locally valid
yet unsupported as a full chain. An interval can be large yet unable to
change the preferred option. A representation edit can invalidate one
receipt while leaving another transportable.

Our thesis is therefore not ``typed state is good.'' It is:

\begin{quote}
\textbf{A partial-knowledge state is useful only relative to the
responsibility that consumes it; at matched visible information, erasing
the responsibility-relevant type/scope can change the decision.}
\end{quote}

This is tested through six exact-synthetic mechanism worlds with
prospectively frozen protocols, matched-information controls,
strongest-donor first right of refusal and hostile validity regimes.

\hypertarget{common-experimental-contract}{%
\subsection{Common experimental
contract}\label{common-experimental-contract}}

\hypertarget{matched-visible-information}{%
\subsubsection{Matched visible
information}\label{matched-visible-information}}

For every primary comparison, all non-oracle arms receive the same
serialized world facts. The paper does not give the candidate privileged
hidden evidence. Instead the candidate and control consume the same
information differently.

Examples: - N4-A: same graph/unknown edges/type labels/rates; the
isolation changes how type-conditioned priors are used; - N4-B: same
failure receipt and observed coordinate changes; the isolation removes
the receipt's dependency scope; - N4-C: same intervals/verification
budget; the isolation changes which uncertainty is considered
decision-relevant; - N4-D: same serialized transport chain; the
isolation changes whether support is checked over the whole chain or
only a local summary/hop; - N4-E: same unknowns/entropy; the isolation
changes whether probes are scored for decision impact; - N4-F3: same
edit history/receipts; the isolation changes whether typed invalidation
governs reuse/remint.

The reviewer-facing parity map is stored in
\texttt{INFORMATION\_PARITY\_AND\_ARTIFACT\_MAP\_V2.md}. Final
submission should convert it to a machine-readable manifest.

\hypertarget{first-right-of-refusal}{%
\subsubsection{First right of refusal}\label{first-right-of-refusal}}

Each study registers strong alternatives before outcome access. A donor
tie or win is an admissible endpoint. The candidate does not receive
novelty credit for a policy or rule already closed by a stronger donor
given the same facts.

\hypertarget{hostile-validity-regimes}{%
\subsubsection{Hostile validity
regimes}\label{hostile-validity-regimes}}

Each positive world contains a construction designed to punish a
specific shortcut. If the hostile control does not bite, the protocol
declares the world invalid rather than counting a candidate advantage.
This is how the study separates ``candidate did well'' from ``the
intended mechanism was actually exercised.''

\hypertarget{evidence-class}{%
\subsubsection{Evidence class}\label{evidence-class}}

The worlds are deterministic exact-synthetic stress tests. Episode
counts are generated units under named frozen generators, not samples
from a population of real scientific domains. Replay establishes
determinism, not independent scientific replication.

\hypertarget{typed-priors-which-unknown-is-worth-probing}{%
\subsection{Typed priors: which unknown is worth
probing?}\label{typed-priors-which-unknown-is-worth-probing}}

N4-A uses a layered interface graph where unknown edge feasibility
depends on an exposed interface type. The candidate and isolation
control use the same myopic VOI planner; the difference is whether the
frozen type-conditioned rates are retained or flattened to a uniform
prior.

Across 300 paired generated episodes, full-oracle mean utility is 4.612.
\texttt{ORION\_TYPED\_VOI} reaches 3.291 with 1.39 probes/episode,
versus 2.180 for the identical uniform-prior planner and 0.358 for exact
optimization restricted to the known subgraph. Blind optimistic
commitment is punished to -13.619, satisfying the hostile validity
condition.

The result does not say an agent can learn the correct type in the wild.
The type/rate facts are supplied by the frozen world. The result says
that \textbf{given those same visible facts}, flattening the registered
type distinction changes the value assigned to unknowns and reduces
utility in this construction.

\hypertarget{scope-invalidation-when-does-an-old-failure-still-apply}{%
\subsection{Scope invalidation: when does an old failure still
apply?}\label{scope-invalidation-when-does-an-old-failure-still-apply}}

N4-B stores a failure receipt together with the exact context
coordinates it depends on. The world also contains an irrelevant
\texttt{NOISE} coordinate that changes frequently. A receipted failure
becomes potentially stale only when a coordinate \emph{inside its
recorded scope} changes.

This is not the same problem as STALE-style memory invalidation. No new
observation semantically contradicts the old failure. The question is
whether the \textbf{conditions under which the failure was established
still match the current task state}.

Policies include \texttt{ORION\_SCOPED\_REOPEN}, \texttt{NEVER\_REOPEN},
\texttt{ALWAYS\_REOPEN} and \texttt{UNSCOPED\_CHANGE\_REOPEN}, which
reopens after any coordinate change including \texttt{NOISE}.

Pooled mean utility is 3.199 for scope-bound reopening, 2.782 for never
reopen, -7.813 for unscoped reopening and -9.225 for always reopen. In
\texttt{STALE\_MATTERS}, scope-bound reopening reaches 2.870 versus
2.096 for never reopen. In \texttt{REOPEN\_WASTEFUL}, scope-bound
reopening remains 3.528 versus 3.468 for never reopen, while always
reopen falls to -13.406 and unscoped reopening to -11.522.

Thus a change event is not automatically a revocation event. The receipt
needs an explicit dependency scope if the downstream responsibility is
``does this old failure still apply?''

\hypertarget{interval-verification-which-uncertainty-can-change-the-choice}{%
\subsection{Interval verification: which uncertainty can change the
choice?}\label{interval-verification-which-uncertainty-can-change-the-choice}}

N4-C gives every arm the same interval-valued cost/error state and the
same budget of four exact edge verifications. The candidate identifies
interval-dominance-surviving paths and prioritizes uncertainty that
participates in unresolved comparisons among them; a matched control
spends the same verification budget randomly.

Across 400 paired generated episodes, mean scalarized regret is 0.1096
for targeted verification and 0.2518 for random verification. The ratio
is 2.3×, but the absolute regrets are the primary effect statement.
Midpoint optimization yields 0.2621, robust worst-case 0.7755 and
best-case 1.2679. The candidate has zero regret in 76.5\% of episodes.

The result is not a universal active-learning theorem. It isolates a
representation/decision relationship: \textbf{uncertainty matters when
it can change the decision boundary, not merely when its interval is
wide.}

\hypertarget{transport-local-validity-is-not-chain-authority}{%
\subsection{Transport: local validity is not chain
authority}\label{transport-local-validity-is-not-chain-authority}}

N4-D contains 200 honest and 200 hostile transport chains in a
constructed finite battery. Hostile classes include missing interior
receipts, spoofed summary tiers and deep splices whose final hop looks
locally legal while an interior input/output identity breaks the chain.

All checkers receive the same serialized chains. Full-chain checking
detects all 200 registered hostile chains and rejects none of the 200
honest chains in this model. Label matching and summary-tier checking
have zero recall; last-hop checking reaches 0.085 overall and zero on
all 68 deep splices.

This is not a cryptographic/security guarantee against real adversaries.
The synthetic model assumes the checker receives the declared receipt
objects. The conclusion is structural: \textbf{authority to
reuse/transport a result can depend on the support path, not only on the
last local label.}

\hypertarget{experiment-selection-entropy-and-decision-relevance-differ}{%
\subsection{Experiment selection: entropy and decision relevance
differ}\label{experiment-selection-entropy-and-decision-relevance-differ}}

N4-E uses a shared stopping/commit rule and adds high-entropy decoy
facts that cannot change the decision. Pure information gain is
therefore tempted by uncertainty that is scientifically irrelevant to
the registered responsibility.

Decision-coupled probing reaches mean utility 9.266, compared with 7.121
for pure information gain, 8.075 for cheapest-first, 7.568 random and
8.989 for the declared deterministic LLM-proxy heuristic. All arms reach
final commit accuracy 1.0; the difference is experiment-selection
efficiency. Pure information gain spends 36.6\% of probes on decoys
while the decision-coupled selector spends zero.

Again the claim is bounded: this world demonstrates that information
value and decision value can diverge even under identical visible
uncertainty.

\hypertarget{representation-edits-reuse-remint-or-re-derive}{%
\subsection{Representation edits: reuse, remint or
re-derive?}\label{representation-edits-reuse-remint-or-re-derive}}

N4-F3 studies whether receipts remain transportable across
representation edits. A typed invalidation rule decides which state can
be reused or reminted. Controls include matched-budget re-derivation and
naive carry-forward.

In \texttt{MIXED\_TRANSPORT}, typed remint/transport reaches 9.421
versus 7.157 for matched re-derivation and -7.821 for naive
carry-forward. Pooled values are 7.286, 6.439 and -3.976 respectively,
with zero registered invalidation mismatches across 14,400 receipts.

The first-right-of-refusal control is decisive: in
\texttt{REMINT\_UNNECESSARY}, all four arms tie exactly at
11.809659685355605 and the candidate spends zero remints. The paper
therefore does not claim that typed reminting always helps; its value is
conditional on the state distinction actually being needed.

\hypertarget{two-negativedonor-outcomes}{%
\subsection{Two negative/donor
outcomes}\label{two-negativedonor-outcomes}}

\hypertarget{typed-failure-state-without-allocation-policy-novelty}{%
\subsubsection{Typed failure state without allocation-policy
novelty}\label{typed-failure-state-without-allocation-policy-novelty}}

N1-C gives the typed state a bounded solve-rate advantage over an
unscoped ablation (+0.0271 paired, bootstrap 95\% interval
\texttt{{[}0.0248,0.02955{]}}) and removes false escalation. But an
ideal VOI donor given the same typed facts \textbf{exactly matches the
candidate allocation policy} at solve rate 0.9866. The allowed claim is
state-value, not policy novelty.

\hypertarget{crossover-predictor-absorbed-on-the-original-world}{%
\subsubsection{Crossover predictor absorbed on the original
world}\label{crossover-predictor-absorbed-on-the-original-world}}

N2-F5B compares a candidate predictor with a model-selection donor. They
tie at 0.9948 in the original world. The candidate retains only a
bounded misspecification-robustness edge in a separately frozen
misspecified world. The original residual is therefore donor-absorbed.

These negatives are central because they show that the methodology can
return ``donor sufficient'' rather than forcing every study into a
typed-state win.

\hypertarget{relation-to-provenancecontext-governance-and-orion-23}{%
\subsection{Relation to provenance/context governance and
ORION-23}\label{relation-to-provenancecontext-governance-and-orion-23}}

ORION-04 assumes that relevant context can be identified, serialized and
compared under exact synthetic rules. \textbf{ContextNest} and
database/workflow provenance research own substantial infrastructure for
version identity, provenance, deterministic selection and
reconstructable history. ORION-04 does not replace that layer.

Its question sits downstream:

\begin{quote}
\textbf{once trustworthy state is available, which distinctions in that
state are required by the next responsibility?}
\end{quote}

The paper is also narrower than ORION-23. ORION-23 develops a general
theory/contract of \textbf{responsibility-scoped sufficiency and
recovery}. ORION-04 supplies bounded exact-synthetic mechanism-isolation
evidence for several decisions that motivate such a theory. ORION-04
therefore does not claim the general principle ``sufficiency is
responsibility-scoped'' as its unique theoretical novelty; it provides
controlled experimental support/examples for specific
responsibility/state distinctions.

\hypertarget{cross-study-synthesis}{%
\subsection{Cross-study synthesis}\label{cross-study-synthesis}}

The six primary worlds use different algorithms but share the same
logic.

\begin{longtable}[]{@{}
  >{\raggedright\arraybackslash}p{(\columnwidth - 4\tabcolsep) * \real{0.3333}}
  >{\raggedright\arraybackslash}p{(\columnwidth - 4\tabcolsep) * \real{0.3333}}
  >{\raggedright\arraybackslash}p{(\columnwidth - 4\tabcolsep) * \real{0.3333}}@{}}
\toprule\noalign{}
\begin{minipage}[b]{\linewidth}\raggedright
Downstream responsibility
\end{minipage} & \begin{minipage}[b]{\linewidth}\raggedright
Required state distinction
\end{minipage} & \begin{minipage}[b]{\linewidth}\raggedright
Matched shortcut that fails in the frozen world
\end{minipage} \\
\midrule\noalign{}
\endhead
\bottomrule\noalign{}
\endlastfoot
choose which unknown to probe & interface/type-conditioned feasibility &
flatten all unknowns to one prior \\
decide whether old failure still applies & dependency scope & reopen on
any change / never reopen \\
spend verification budget & uncertainty that can change choice & verify
arbitrary/high-width uncertainty \\
accept transported evidence & full-chain support/continuity & trust
final label/summary/last hop \\
choose next experiment & decision relevance & maximize entropy alone \\
reuse state after representation edit & transport/invalidation type &
rederive everything or carry everything forward \\
\end{longtable}

The commonality is not a universal type system. It is a design
criterion: \textbf{the state representation should retain exactly the
distinctions needed by the downstream responsibility, while
first-right-of-refusal controls test whether a simpler donor already
suffices.}

\hypertarget{reproducibility-and-statistical-reporting}{%
\subsection{Reproducibility and statistical
reporting}\label{reproducibility-and-statistical-reporting}}

Every primary world is generated from a frozen protocol and
deterministic code. The final package should bind one machine-readable
manifest per study with:

\begin{itemize}
\tightlist
\item
  generator/seed;
\item
  exact serialized information visible to each arm;
\item
  policy difference;
\item
  independent generated unit/episode count;
\item
  hostile-control gate;
\item
  source result digest;
\item
  replay command.
\end{itemize}

Do not pool all episodes/chains across the six worlds into one p value
or meta-effect. They are different constructed mechanisms with different
outcomes/scales. For N4-D, 200 hostile/200 honest chains are an exact
finite battery rather than IID draws from a real adversary population.
For N1-C, the registered paired bootstrap interval must preserve the
protocol's actual independent unit.

The repository is publicly inspectable. A final article may describe
code/data as open/reusable only after an authorized licence and
permanent archive are actually established.

\hypertarget{limitations}{%
\subsection{Limitations}\label{limitations}}

\textbf{Exact-synthetic worlds.} No primary result is a real-agent,
real-laboratory or deployment study.

\textbf{Rules/types are supplied.} The experiments generally assume the
relevant type/scope facts are present in the state. They do not show
that an agent can learn correct responsibility types automatically.

\textbf{Current donors are broad.} STALE, ContextNest, VOI, active
learning and provenance already own major primitives. ORION-04's
residual is the controlled matched-information mechanism evidence.

\textbf{Constructed hostile regimes.} Hostile controls are intentionally
designed to exercise the target distinction; they are not prevalence
estimates.

\textbf{No cryptographic claim.} Hash/receipt objects in N4-D are
synthetic provenance mechanisms, not proof of security against real
attackers.

\textbf{ORION-23 ownership.} General responsibility-scoped
sufficiency/recovery theory belongs to ORION-23; ORION-04 is bounded
empirical/mechanism evidence.

\textbf{No single universal mechanism.} The six worlds need different
type/scope coordinates because the downstream responsibilities differ.

\hypertarget{discussion}{%
\subsection{Discussion}\label{discussion}}

A useful state representation is not defined only by how much
information it stores. It is defined by which distinctions survive long
enough to matter for the next decision.

ORION-04's strongest examples are cases where a superficially reasonable
scalarization loses that structure. ``Changed'' is weaker than ``a
coordinate inside the receipt's dependency scope changed.''
``Uncertain'' is weaker than ``uncertainty can reverse the decision.''
``Last hop valid'' is weaker than ``the full support chain remains
valid.'' ``High entropy'' is weaker than ``worth probing for the current
choice.''

This interpretation complements rather than competes with current
agent-memory and context-governance work. Those systems help maintain
trustworthy, current state. ORION-04 tests the next layer: \textbf{how
the responsibility consuming that state determines which distinctions
must remain explicit.}

The first-right-of-refusal negatives matter just as much. When an ideal
VOI donor already has the typed facts, the candidate allocation policy
earns no novelty. When model selection closes an original prediction
residual, the paper retains the donor tie. This is evidence that the
experimental discipline is capable of subtracting mechanisms rather than
only accumulating wins.

\hypertarget{conclusion}{%
\subsection{Conclusion}\label{conclusion}}

Across six frozen exact-synthetic research-decision worlds, explicit
type/scope is load-bearing precisely when the registered downstream
responsibility depends on a distinction that a matched shortcut erases.
The result is not a new generic memory, provenance, VOI or
active-learning method. It is a bounded mechanism finding about
\textbf{state consumption under matched information}, with donor ties
and hostile controls preserved as first-class evidence.

The practical lesson is correspondingly narrow: before compressing a
research state into one confidence, validity flag, relevance score or
``changed/not changed'' bit, ask which responsibility must consume it
next---and whether the discarded distinction can change that decision.

\section*{Data and Code Availability}
All primary study worlds in this manuscript are synthetic and generated deterministically
from frozen protocols and repository code. Result receipts for the six primary studies and
two donor/negative controls are bound to the publication evidence cut. The source repository
is publicly inspectable. A permanent archival identifier and explicit reuse licence will be
inserted in the final submission only after the corresponding release is actually deposited
and authorized; no protected or private human dataset is used in this study.

\bibliographystyle{tmlr}
\bibliography{references}

\end{document}
